\documentclass[a4paper,11pt]{report}

% ===================== LANGUE & POLICES =====================
\usepackage[utf8]{inputenc}
\usepackage[T1]{fontenc}
\usepackage[french]{babel}
\usepackage{lmodern}
\usepackage[margin=2.5cm]{geometry}

% ===================== MISE EN PAGE & GRAPHISME =====================
\usepackage{graphicx}
\usepackage{float}
\usepackage{booktabs}
\usepackage{array}
\usepackage{caption}
\usepackage{subcaption}
\usepackage[table]{xcolor}
\usepackage{amsmath}
\usepackage{siunitx}
\usepackage{fancyhdr}
\usepackage{titlesec}
\usepackage{enumitem}
\usepackage{url}
\usepackage{tikz}
\usetikzlibrary{positioning, arrows.meta}
\usepackage{fontawesome5}
\usepackage{listings}
\usepackage[skins,breakable]{tcolorbox}
\tcbuselibrary{listings}
\usepackage[colorlinks=true,linkcolor=primary,citecolor=primary,urlcolor=secondary]{hyperref}

% ===================== PALETTE DE COULEURS =====================
\definecolor{primary}{HTML}{13315C}   % bleu nuit - titres, cadres
\definecolor{secondary}{HTML}{2C7DA0} % bleu clair - accents, liens
\definecolor{accent}{HTML}{D9822B}    % orange - avertissements
\definecolor{success}{HTML}{2E7D32}   % vert - validations
\definecolor{codebg}{HTML}{F5F7FA}    % fond de code
\definecolor{codegreen}{HTML}{2E7D32}
\definecolor{codegray}{HTML}{6B7280}
\definecolor{codepurple}{HTML}{8B5CF6}

% ===================== EN-TÊTE / PIED DE PAGE =====================
\setlength{\headheight}{22pt}
\pagestyle{fancy}
\fancyhf{}
\renewcommand{\headrulewidth}{0.8pt}
\renewcommand{\headrule}{\hbox to\headwidth{\color{secondary}\leaders\hrule height \headrulewidth\hfill}}
\fancyhead[L]{\small\color{primary}\textsc{Infrastructure de capteurs LoRa/MQTT}}
\fancyhead[R]{\small\color{primary}\leftmark}
\renewcommand{\chaptermark}[1]{\markboth{#1}{}}
\fancyfoot[C]{\small\color{codegray}\thepage}
\fancyfoot[R]{\small\color{codegray}\textsc{INSA - STI 4A}}

% ===================== TITRES DE CHAPITRES / SECTIONS =====================
\titleformat{\chapter}[display]
  {\normalfont\bfseries\color{primary}}
  {\textsc{\large Chapitre \thechapter}}
  {10pt}
  {\Huge\bfseries}
  [{\vspace{6pt}\noindent\textcolor{secondary}{\rule{\linewidth}{1pt}}}]
\titlespacing*{\chapter}{0pt}{0pt}{25pt}

\titleformat{\section}
  {\normalfont\Large\bfseries\color{primary}}
  {\textcolor{secondary}{\rule[-2pt]{4pt}{14pt}}\hspace{6pt}\thesection}
  {8pt}{}
\titlespacing*{\section}{0pt}{16pt}{8pt}

\titleformat{\subsection}
  {\normalfont\large\bfseries\color{secondary}}
  {\thesubsection}{6pt}{}
\titlespacing*{\subsection}{0pt}{12pt}{6pt}

% ===================== STYLE DES TABLEAUX =====================
\renewcommand{\arraystretch}{1.25}

% ===================== STYLES DE CODE (listings) =====================
\lstset{
  literate=
    {é}{{\'e}}1 {è}{{\`e}}1 {à}{{\`a}}1 {ù}{{\`u}}1
    {â}{{\^a}}1 {ê}{{\^e}}1 {î}{{\^i}}1 {ô}{{\^o}}1 {û}{{\^u}}1
    {ç}{{\c{c}}}1 {œ}{{\oe}}1
    {É}{{\'E}}1 {È}{{\`E}}1 {À}{{\`A}}1 {Ç}{{\c{C}}}1
}

\lstdefinestyle{pystyle}{
    language=Python,
    commentstyle=\color{codegreen}\itshape,
    keywordstyle=\color{primary}\bfseries,
    stringstyle=\color{codepurple},
    numberstyle=\tiny\color{codegray},
    basicstyle=\ttfamily\small,
    numbers=left,
    numbersep=8pt,
    breaklines=true,
    tabsize=4,
    showstringspaces=false,
}

\lstdefinestyle{shellstyle}{
    basicstyle=\ttfamily\small\color{white},
    breaklines=true,
    showstringspaces=false,
}

% ===================== BOÎTES PERSONNALISÉES =====================
% Bloc de code "fichier" avec bandeau de titre
\newtcblisting{pycode}[1]{
    listing only,
    listing options={style=pystyle},
    breakable, enhanced,
    colback=codebg, colframe=primary,
    coltitle=white, fonttitle=\bfseries\ttfamily\small,
    title={\faFileCode\ #1},
    boxrule=0.6pt, arc=2pt, top=2mm, bottom=1mm, left=1mm, right=1mm,
    attach boxed title to top left={yshift=-2mm, xshift=2mm},
    boxed title style={colback=primary, boxrule=0pt, arc=2pt},
}

% Bloc "terminal" pour les commandes shell
\newtcblisting{shellcode}[1]{
    listing only,
    listing options={style=shellstyle},
    breakable, enhanced,
    colback=black!88, colframe=black!60,
    coltitle=white, fonttitle=\bfseries\ttfamily\small,
    title={\faTerminal\ #1},
    boxrule=0.6pt, arc=2pt, top=2mm, bottom=1mm, left=1mm, right=1mm,
}

% Encadré d'explication pédagogique
\newtcolorbox{explication}[1][Explication]{
    breakable, enhanced,
    colback=secondary!6, colframe=secondary,
    coltitle=white, fonttitle=\bfseries,
    title={\faLightbulb\ #1},
    boxrule=0.6pt, arc=2pt, left=2mm, right=2mm,
}

% Encadré de remarque / point d'attention
\newtcolorbox{remarque}[1][À noter]{
    breakable, enhanced,
    colback=accent!8, colframe=accent,
    coltitle=white, fonttitle=\bfseries,
    title={\faExclamationTriangle\ #1},
    boxrule=0.6pt, arc=2pt, left=2mm, right=2mm,
}

% Encadré de validation / résultat positif
\newtcolorbox{validation}[1][Résultat]{
    breakable, enhanced,
    colback=success!6, colframe=success,
    coltitle=white, fonttitle=\bfseries,
    title={\faCheckCircle\ #1},
    boxrule=0.6pt, arc=2pt, left=2mm, right=2mm,
}

% ===================== INFOS DU PROJET =====================
\newcommand{\projecttitle}{Infrastructure de capteurs : broker MQTT \& transmission LoRa sécurisée}
\newcommand{\authorone}{BARRY Boubacar Sidy}
\newcommand{\authortwo}{Nhat Lam NGUYEN}
\newcommand{\formation}{INSA - STI 4A - RiSF}
\newcommand{\anneeuniv}{2025 - 2026}

\begin{document}

% =====================================================
% PAGE DE GARDE
% =====================================================
\begin{titlepage}
    \centering
    {\color{secondary}\rule{\textwidth}{2pt}}

    \vspace{0.6cm}
    {\scshape\Large\color{primary} \formation \par}
    \vspace{0.2cm}
    {\large\color{codegray} Année universitaire \anneeuniv \par}

    \vspace{2cm}

    \begin{tcolorbox}[
        enhanced, sharp corners, boxrule=0pt,
        colback=primary, coltext=white,
        width=\textwidth, halign=center,
        drop shadow={black!40},
        top=10mm, bottom=10mm
    ]
        {\large\scshape Rapport de projet\par}
        \vspace{0.5cm}
        {\Huge\bfseries \projecttitle \par}
    \end{tcolorbox}

    \vspace{1.8cm}

    \begin{tcolorbox}[
        enhanced, colback=white, colframe=secondary,
        width=0.7\textwidth, halign=center, arc=3pt
    ]
        {\Large\bfseries \authorone \par}
        \vspace{2pt}
        {\Large\bfseries \authortwo \par}
    \end{tcolorbox}

    \vfill

    {\large Sujet : \emph{Fiche INSA RiSF LoRa - Question 6} \par}
    {\large \today \par}

    \vspace{0.4cm}
    {\color{secondary}\rule{\textwidth}{2pt}}
\end{titlepage}

% =====================================================
% RÉSUMÉ
% =====================================================
\renewcommand{\abstractname}{Résumé}
\begin{abstract}
\noindent
Ce rapport présente la réalisation d'une infrastructure de capteurs sans fil combinant LoRa et MQTT.
Un premier ESP32 joue le rôle de capteur LoRa et publie périodiquement des mesures chiffrées.
Un second ESP32 agit comme passerelle : il reçoit les trames LoRa, vérifie leur authenticité,
déchiffre les données et les republie vers un broker MQTT. Enfin, un client Python exécuté sur Linux
s'abonne au broker, vérifie les messages reçus et permet de publier des trames de test ou de mettre
à jour les clés de sécurité. Le protocole applique un chiffrement AES-CBC et une authentification par
HMAC-SHA1 afin de protéger la confidentialité, l'intégrité et l'origine des données.
\end{abstract}

\tableofcontents
\listoffigures
\listoftables

% =====================================================
\chapter{Introduction}
\label{chap:intro}

\section{Contexte}
Le projet s'inscrit dans le cadre du cours Réseaux Sans Fil de l'INSA. L'objectif est de mettre en
œuvre une chaîne complète de communication IoT : génération d'une mesure par un capteur, transmission
radio longue portée via LoRa, réception par une passerelle, puis diffusion des données vers un broker
MQTT. Le broker permet ensuite à un client distant de consulter les mesures et de vérifier que les
messages reçus sont bien authentiques et intègres.

\section{Objectifs}
Les objectifs principaux sont les suivants :
\begin{itemize}
    \item transmettre des données capteur entre deux ESP32 à l'aide de LoRa ;
    \item republier les données reçues vers un broker MQTT ;
    \item chiffrer les messages avec AES-CBC ;
    \item ajouter une authentification par HMAC-SHA1 ;
    \item vérifier côté passerelle puis côté client Linux que les messages n'ont pas été modifiés ;
    \item permettre une mise à jour sécurisée des clés AES et HMAC depuis le client Python.
\end{itemize}

\section{Organisation du rapport}
Le rapport décrit d'abord l'architecture générale, puis le fonctionnement détaillé - fichier par
fichier - du capteur, de la passerelle et du client Linux. Les mécanismes de sécurité sont ensuite
analysés avant de présenter les tests, les difficultés rencontrées et les améliorations possibles.

% =====================================================
\chapter{Architecture générale de la solution}
\label{chap:archi}

\section{Vue d'ensemble}
L'architecture comprend quatre éléments principaux :
\begin{enumerate}
    \item un ESP32 \textbf{capteur} qui génère une mesure et l'envoie en LoRa ;
    \item un ESP32 \textbf{passerelle} qui reçoit la trame LoRa, la vérifie et la republie ;
    \item un \textbf{broker MQTT} distant (\texttt{serveur.iot.com} ou \texttt{p-fb.net}) ;
    \item un \textbf{client Python} sous Linux qui s'abonne au topic MQTT et vérifie les messages.
\end{enumerate}

\begin{figure}[H]
    \centering
    \begin{tikzpicture}[
        node distance=1.4cm,
        box/.style={draw=primary, fill=primary!6, rounded corners=2pt, minimum height=1.1cm, minimum width=2.6cm, align=center, font=\small\bfseries\color{primary}, thick},
        arrow/.style={-{Stealth[length=3mm]}, secondary, thick}
    ]
        \node[box] (capteur) {Capteur\\ESP32 \#1};
        \node[box, right=of capteur] (passerelle) {Passerelle\\ESP32 \#2};
        \node[box, right=of passerelle] (broker) {Broker\\MQTT};
        \node[box, right=of broker] (client) {Client\\Linux};

        \draw[arrow] (capteur) -- (passerelle) node[midway, above, font=\scriptsize\color{secondary}] {LoRa};
        \draw[arrow] (passerelle) -- (broker) node[midway, above, font=\scriptsize\color{secondary}] {MQTT};
        \draw[arrow] (broker) -- (client) node[midway, above, font=\scriptsize\color{secondary}] {MQTT};
        \draw[arrow, dashed] (client.south) -- ++(0,-0.8) -| (passerelle.south) node[midway, below, font=\scriptsize\color{accent}] {commande de mise à jour des clés};
    \end{tikzpicture}
    \caption{Architecture générale de l'infrastructure}
    \label{fig:architecture}
\end{figure}

\section{Matériel et environnement logiciel}
\begin{itemize}
    \item 2 cartes ESP32 ;
    \item 2 modules LoRa SX127x ;
    \item MicroPython sur ESP32 ;
    \item Broker MQTT : \texttt{serveur.iot.com} ou \texttt{p-fb.net}, port \texttt{1883} ;
    \item Client Linux : Python 3 avec \texttt{paho-mqtt} et \texttt{cryptography}.
\end{itemize}

\section{Format des trames échangées}
Les messages LoRa sont encodés en base64. Le format logique est :
\[
\texttt{node\_id | nonce + ciphertext\_AES | HMAC}.
\]
Le HMAC est calculé sur :
\[
\texttt{INSARISF1 | node\_id | nonce + ciphertext\_AES}.
\]
La clé AES fait 16 octets et le mode utilisé est AES-CBC. Le nonce correspond à l'IV AES de 16 octets.
Le HMAC utilise SHA-1 et produit un tag de 20 octets.

\begin{table}[H]
    \centering
    \caption{Format de message applicatif}
    \label{tab:trame}
    \rowcolors{2}{primary!6}{white}
    \begin{tabular}{lll}
        \toprule
        \rowcolor{primary} \textcolor{white}{Champ} & \textcolor{white}{Taille} & \textcolor{white}{Description} \\
        \texttt{node\_id} & variable & Identifiant du nœud émetteur \\
        \texttt{nonce} & 16 octets & IV AES aléatoire \\
        \texttt{ciphertext} & multiple de 16 octets & Données chiffrées AES-CBC \\
        \texttt{HMAC-SHA1} & 20 octets & Intégrité et authenticité \\
        Encodage final & base64 & Transport LoRa et MQTT \\
        \bottomrule
    \end{tabular}
\end{table}

% =====================================================
\chapter{Module capteur (ESP32 \#1)}
\label{chap:capteur}

Ce chapitre détaille le code embarqué qui s'exécute sur le premier ESP32, jouant le rôle de
\textbf{capteur LoRa}. Son unique mission est de produire une mesure, de la sécuriser, puis de
l'émettre en LoRa - il ne se connecte ni au Wi-Fi, ni au broker MQTT.

\section{Fichier de configuration \texttt{config.py}}

\begin{pycode}{Projet/firmware/capteur/config.py}
from machine import Pin, SoftSPI

WIFI_SSID = 'INSA_RiSF'
WIFI_PASSWORD = 'insa2026'

MQTT_SERVER = 'serveur.iot.com'
MQTT_PORT = 1883
MQTT_USER = 'esp32'
MQTT_PASSWORD = 'esp32'
MQTT_DATA_TOPIC = b'LamBouba'
MQTT_COMMAND_TOPIC = b'echange/command/passerelle'

NODE_ID = 'capteur1'
AES_KEY = b'0123456789abcdef'
HMAC_KEY = b'abcd'

device_config = {
    'miso': 19,
    'mosi': 27,
    'ss': 18,
    'sck': 5,
    'dio_0': 26,
    'reset': 14,
    'led': 25,
}

lora_parameters = {
    'frequency': 868E6,
    'tx_power_level': 2,
    'signal_bandwidth': 125E3,
    'spreading_factor': 11,
    'coding_rate': 5,
    'preamble_length': 8,
    'implicit_header': True,
    'sync_word': 0x12,
    'enable_CRC': True,
    'invert_IQ': False,
}
\end{pycode}

\begin{explication}[Un gabarit commun pour les deux nœuds]
Le fichier \texttt{config.py} du capteur partage la même structure que celui de la passerelle.
Cette conception facilite la maintenance : modifier une constante (fréquence LoRa, clé AES) dans
un gabarit unique suffit pour les deux firmwares. Les constantes réseau (\texttt{WIFI\_SSID},
\texttt{MQTT\_*}) sont définies ici et activement utilisées par la passerelle.
\end{explication}

\subsection{Identité et secrets partagés}
\begin{itemize}
    \item \texttt{NODE\_ID} identifie le capteur dans les trames applicatives (\texttt{capteur1}) ;
    il doit être unique sur le réseau LoRa.
    \item \texttt{AES\_KEY} (16 octets) est la clé symétrique utilisée pour chiffrer la charge utile en
    AES-CBC. Elle doit être strictement identique sur le capteur, la passerelle et le client Linux.
    \item \texttt{HMAC\_KEY} sert à authentifier les messages \emph{et} les commandes de mise à jour des
    clés. Elle doit être identique sur le capteur, la passerelle et le client Linux, et peut être mise
    à jour dynamiquement via le protocole de rotation des clés (voir chapitre~\ref{chap:securite}).
\end{itemize}

\subsection{Brochage \texttt{device\_config}}
Ce dictionnaire décrit le câblage SPI entre l'ESP32 et le module radio SX127x :
\begin{itemize}
    \item \texttt{sck}, \texttt{mosi}, \texttt{miso} : les trois lignes de l'horloge et des données du
    bus SPI logiciel (\texttt{SoftSPI}) ;
    \item \texttt{ss} : la ligne de sélection de puce (\emph{chip select} / NSS) du module LoRa ;
    \item \texttt{dio\_0} : la broche d'interruption du SX127x, utilisée par le pilote pour détecter la
    fin d'une réception ou d'une émission sans avoir à interroger le registre en continu ;
    \item \texttt{reset} : la broche de réinitialisation matérielle du module radio ;
    \item \texttt{led} : une broche libre pour une éventuelle LED de diagnostic.
\end{itemize}

\subsection{Paramètres radio \texttt{lora\_parameters}}
Ces paramètres définissent la couche physique LoRa et \textbf{doivent être identiques} sur le capteur
et sur la passerelle, sous peine de ne jamais pouvoir communiquer :
\begin{itemize}
    \item \texttt{frequency} : \SI{868}{\mega\hertz}, bande ISM européenne autorisée pour le LoRa ;
    \item \texttt{tx\_power\_level} : puissance d'émission demandée à l'amplificateur du SX127x ;
    \item \texttt{signal\_bandwidth} : largeur de bande du signal (\SI{125}{\kilo\hertz}), qui conditionne
    le débit symbole ;
    \item \texttt{spreading\_factor} : facteur d'étalement (ici SF11). Plus il est élevé, plus la portée
    et la résistance au bruit augmentent, mais plus le temps d'émission (\emph{airtime}) d'un paquet
    s'allonge ;
    \item \texttt{coding\_rate} : taux de codage correcteur d'erreurs (ici 4/5), qui ajoute de la
    redondance pour corriger les erreurs de transmission ;
    \item \texttt{preamble\_length} : nombre de symboles de préambule envoyés avant les données utiles,
    nécessaires à la synchronisation du récepteur ;
    \item \texttt{implicit\_header} : à \texttt{True}, l'en-tête LoRa (longueur, coding rate, présence de
    CRC) n'est pas transmis dans la trame radio - émetteur et récepteur doivent donc partager
    \emph{a priori} exactement la même configuration ;
    \item \texttt{sync\_word} : mot de synchronisation (\texttt{0x12}) qui permet de filtrer les trames
    appartenant à un autre réseau LoRa utilisant la même fréquence ;
    \item \texttt{enable\_CRC} : active un contrôle d'intégrité au niveau de la couche physique, en plus
    du HMAC applicatif ;
    \item \texttt{invert\_IQ} : inversion I/Q, désactivée ici (utile uniquement pour certains modes
    spécifiques comme le LoRaWAN \emph{downlink}).
\end{itemize}

\section{Programme principal \texttt{main.py}}

\begin{pycode}{Projet/firmware/capteur/main.py}
import os
import sys
os.chdir("/capteur")
from machine import Pin, SoftSPI
from time import sleep, ticks_ms, ticks_diff
import machine
import ujson

from config import *
from security import *
from sx127x import SX127x
from oled import *


def init_lora():
    device_spi = SoftSPI(
        baudrate=10000000,
        polarity=0,
        phase=0,
        bits=8,
        firstbit=SoftSPI.MSB,
        sck=Pin(device_config['sck'], Pin.OUT, Pin.PULL_DOWN),
        mosi=Pin(device_config['mosi'], Pin.OUT, Pin.PULL_UP),
        miso=Pin(device_config['miso'], Pin.IN, Pin.PULL_UP)
    )
    return SX127x(
        device_spi,
        pins=device_config,
        parameters=lora_parameters
    )


def build_payload(counter):
    value = 20.0 + counter * 0.5
    return ujson.dumps({
        'node': NODE_ID,
        'type': 'sensor',
        'counter': counter,
        'temperature': value,
        'battery': max(0, 90 - counter)
    })


def listen_key_update(lora, oled, screen, timeout_ms=2000):
    global AES_KEY
    global HMAC_KEY

    # Passage explicite en mode reception active
    lora.receive()

    start_time = ticks_ms()
    packet_detected = False

    # Attente active non-bloquante du paquet (evite le sleep brut)
    while ticks_diff(ticks_ms(), start_time) < timeout_ms:
        if lora.received_packet():
            print("Msg recu")
            packet_detected = True
            break
        sleep(0.05)

    if not packet_detected:
        # Forcer le retour au mode veille pour pouvoir emettre au prochain tour
        lora.sleep()
        return False

    # Lecture du paquet de mise a jour des cles
    packet = lora.read_payload()

    # Securite : verification de la signature HMAC du paquet d'update
    ok, new_aes_key, new_hmac_key = verify_key_update(packet, HMAC_KEY)

    if ok:
        AES_KEY = new_aes_key
        HMAC_KEY = new_hmac_key

        screen[0] = 'KEYS OK'
        screen[1] = NODE_ID
        screen[2] = 'updated'
        write_screen(oled, screen)
        lora.sleep()
        return True

    # Paquet recu mais invalide (mauvais HMAC, rejeu...)
    screen[0] = 'KEYS BAD'
    screen[1] = NODE_ID
    screen[2] = 'Rejected!'
    write_screen(oled, screen)
    lora.sleep()
    return False


def run():
    global AES_KEY
    global HMAC_KEY

    lora = init_lora()
    oled, screen = init_oled()
    counter = 0

    screen[0] = 'LoRa capteur'
    screen[1] = NODE_ID
    screen[2] = 'AES+HMAC'
    write_screen(oled, screen)
    sleep(2)

    while True:
        counter += 1
        payload = build_payload(counter)

        # Chiffrement AES-CBC + Signature HMAC-SHA1 + Encodage Base64
        packet = pack_message(NODE_ID, payload.encode(), AES_KEY, HMAC_KEY)

        # Envoi de la trame Base64 via LoRa
        lora.println(packet)
        print("[TX] Paquet envoye : {}".format(packet[:30] + "..."))

        screen[0] = 'TX {}'.format(NODE_ID)
        screen[1] = packet[:16]
        screen[2] = 'counter {}'.format(counter)
        write_screen(oled, screen)

        # Fenetre d'ecoute des cles (6 secondes d'ouverture)
        print("Enter listening")
        listen_key_update(lora, oled, screen, timeout_ms=6000)

if __name__ == '__main__':
    run()
\end{pycode}

\subsection{Imports et organisation}
Le script importe massivement les modules locaux avec \texttt{from ... import *}
(\texttt{config}, \texttt{security}, \texttt{oled}). Cette pratique, déconseillée dans un projet Python
classique, est courante en MicroPython embarqué : elle évite de préfixer chaque constante et limite la
quantité de code source à compiler en mémoire flash. Toutes les constantes de \texttt{config.py} (NODE\_ID,
AES\_KEY, lora\_parameters...) et toutes les fonctions de \texttt{security.py} (\texttt{pack\_message},
\texttt{verify\_key\_update}...) deviennent ainsi directement accessibles dans \texttt{main.py}.

\subsection{\texttt{init\_lora()} - initialisation du module radio}
Le bus SPI n'est pas le bus matériel de l'ESP32 mais un bus \textbf{logiciel} (\texttt{SoftSPI}), ce qui
permet de choisir librement les broches \texttt{sck}/\texttt{mosi}/\texttt{miso} dans
\texttt{device\_config} plutôt que d'être contraint par le brochage SPI matériel par défaut. Le pilote
\texttt{SX127x} est ensuite instancié avec ce bus, le dictionnaire de broches et les paramètres radio :
il configure les registres du module (fréquence, SF, BW, sync word...) lors de sa construction.

\subsection{\texttt{build\_payload()} - génération de la mesure}
Cette fonction simule un capteur physique : la température suit une rampe linéaire
(\texttt{20.0 + counter * 0.5}) et la batterie décroît avec le compteur, sans jamais devenir négative
grâce à \texttt{max(0, ...)}. Le résultat est sérialisé en JSON avec \texttt{ujson.dumps}, l'implémentation
JSON allégée de MicroPython. Dans une version finale, cette fonction serait remplacée par une vraie
lecture de capteur (DHT22, BME280, etc.) sans changer le reste du programme.

\subsection{\texttt{listen\_key\_update()} - fenêtre de réception des clés}
Le capteur est \emph{half-duplex} : il ne peut pas émettre et écouter en même temps. Pour permettre
malgré tout à la passerelle de lui transmettre de nouvelles clés, le capteur ouvre une courte fenêtre
d'écoute après chaque émission :
\begin{enumerate}
    \item \texttt{lora.receive()} arme le module en mode réception ;
    \item une boucle d'attente active interroge \texttt{lora.received\_packet()} toutes les
    \SI{50}{\milli\second}, bornée par \texttt{ticks\_ms()}/\texttt{ticks\_diff()} : cette approche est
    préférée à un simple \texttt{sleep()} bloquant car elle permet de sortir de la boucle dès qu'un paquet
    arrive, sans attendre inutilement tout le \texttt{timeout\_ms} ;
    \item si rien n'est reçu avant l'expiration du délai, le module repasse en veille
    (\texttt{lora.sleep()}) afin de libérer le canal pour la prochaine émission ;
    \item si un paquet est reçu, \texttt{verify\_key\_update()} vérifie son HMAC \emph{avant} d'accepter
    quoi que ce soit ;
    \item en cas de \textbf{succès} : les variables globales \texttt{AES\_KEY} et \texttt{HMAC\_KEY} sont
    réaffectées avec les nouvelles valeurs, puis le capteur envoie un accusé de réception explicite
    \texttt{KEYS\_OK} chiffré et signé avec les \emph{nouvelles} clés via \texttt{pack\_message()} et
    \texttt{lora.println()} ;
    \item en cas d'\textbf{échec} (HMAC invalide, format incorrect) : le capteur répond par un NACK
    \texttt{KEYS\_FAILED} signé avec ses \emph{anciennes} clés toujours actives, ce qui permet à la
    passerelle et au client Linux de détecter le désaccord et de déclencher un rollback automatique ;
    \item dans tous les cas, l'écran OLED affiche le résultat (\texttt{KEYS OK} ou \texttt{KEYS BAD})
    pour un diagnostic visuel immédiat.
\end{enumerate}

\begin{explication}[Accusé de réception avec les nouvelles clés]
Le capteur applique les nouvelles clés \emph{avant} d'envoyer l'ACK, de sorte que le message
\texttt{KEYS\_OK} est lui-même signé avec les nouvelles clés. Le client Linux, qui tente de le déchiffrer
avec les nouvelles clés, peut ainsi confirmer que le capteur les a bien adoptées. Si la rotation a
échoué, le capteur renvoie \texttt{KEYS\_FAILED} signé avec les anciennes clés, ce qui déclenche
le rollback automatique côté client.
\end{explication}

\subsection{\texttt{run()} - boucle principale}
Après une phase d'initialisation (LoRa, écran OLED, affichage de démarrage), le capteur entre dans une
boucle infinie qui, à chaque itération~:
\begin{enumerate}
    \item incrémente le compteur de mesures ;
    \item construit la charge utile JSON avec \texttt{build\_payload()} ;
    \item la sécurise avec \texttt{pack\_message()} (chiffrement AES-CBC puis HMAC-SHA1 puis base64,
    détaillé au chapitre~\ref{chap:client}) ;
    \item émet le paquet en LoRa avec \texttt{lora.println()} ;
    \item met à jour l'écran OLED avec le début du paquet encodé ;
    \item ouvre une fenêtre d'écoute de 6 secondes via \texttt{listen\_key\_update()} avant de recommencer.
\end{enumerate}
Le cycle complet (émission puis écoute) se répète donc indéfiniment, sans délai fixe supplémentaire : la
période d'émission effective est d'environ 6 secondes, dictée par la durée de la fenêtre d'écoute des clés.

% =====================================================
\chapter{Module passerelle (ESP32 \#2)}
\label{chap:passerelle}

La passerelle est le seul nœud qui parle à la fois LoRa et MQTT : elle reçoit les trames radio du
capteur, vérifie leur authenticité, puis les republie sur le broker. Elle joue donc le rôle de
\textbf{client MQTT} pour le reste de l'infrastructure.

\section{Fichier de configuration \texttt{config.py}}

\begin{pycode}{Projet/firmware/passerelle/config.py}
from machine import Pin, SoftSPI

WIFI_SSID = 'INSA_RiSF'
WIFI_PASSWORD = 'insa2026'

MQTT_SERVER = 'serveur.iot.com'
MQTT_PORT = 1883
MQTT_USER = 'esp32'
MQTT_PASSWORD = 'esp32'
MQTT_DATA_TOPIC = b'LamBouba'
MQTT_COMMAND_TOPIC = b'echange/command/passerelle'

NODE_ID = 'passerelle1'
AES_KEY = b'0123456789abcdef'
HMAC_KEY = b'abcd'

device_config = {
    'miso': 19,
    'mosi': 27,
    'ss': 18,
    'sck': 5,
    'dio_0': 26,
    'reset': 14,
    'led': 25,
}

lora_parameters = {
    'frequency': 868E6,
    'tx_power_level': 2,
    'signal_bandwidth': 125E3,
    'spreading_factor': 11,
    'coding_rate': 5,
    'preamble_length': 8,
    'implicit_header': True,
    'sync_word': 0x12,
    'enable_CRC': True,
    'invert_IQ': False,
}
\end{pycode}

Ce fichier reprend exactement la même structure que celui du capteur (chapitre~\ref{chap:capteur}), avec
trois différences importantes :
\begin{itemize}
    \item \texttt{NODE\_ID} vaut \texttt{passerelle1} ;
    \item \texttt{MQTT\_DATA\_TOPIC} (\texttt{LamBouba}) est ici réellement utilisé : c'est le topic sur
    lequel la passerelle publie les données reçues du capteur ;
    \item \texttt{MQTT\_COMMAND\_TOPIC} (\texttt{echange/command/passerelle}) est également utilisé : la
    passerelle s'y abonne pour recevoir les commandes de mise à jour des clés envoyées par le client Linux.
\end{itemize}

\begin{explication}[Synchronisation des paramètres radio]
Les valeurs de \texttt{lora\_parameters} sont \textbf{rigoureusement identiques} entre le capteur et la
passerelle : même fréquence, même facteur d'étalement, même bande passante. C'est cette synchronisation
qui permet aux deux modules SX127x de communiquer, le LoRa exigeant que l'émetteur et le récepteur
partagent exactement la même configuration physique.
\end{explication}

\section{Programme principal \texttt{main.py}}

\begin{pycode}{Projet/firmware/passerelle/main.py}
import os
os.chdir("/passerelle")
from machine import Pin, SoftSPI
from time import sleep
import machine
import network
import ubinascii
import ujson

from config import *
from security import *
from sx127x import SX127x
from oled import *
from umqttsimple import MQTTClient


mqtt_client = None
no_check_mode = False


def init_lora():
    device_spi = SoftSPI(
        baudrate=10000000,
        polarity=0,
        phase=0,
        bits=8,
        firstbit=SoftSPI.MSB,
        sck=Pin(device_config['sck'], Pin.OUT, Pin.PULL_DOWN),
        mosi=Pin(device_config['mosi'], Pin.OUT, Pin.PULL_UP),
        miso=Pin(device_config['miso'], Pin.IN, Pin.PULL_UP)
    )
    modem = SX127x(
        device_spi,
        pins=device_config,
        parameters=lora_parameters
    )
    # Force le modem en mode ecoute LoRa des le demarrage
    modem.receive()
    return modem


def connect_wifi():
    wlan = network.WLAN(network.STA_IF)

    # Nettoyage preventif de l'etat interne de la puce Wi-Fi
    if wlan.active():
        wlan.active(False)
        sleep(0.5)

    wlan.active(True)
    sleep(0.5)

    if wlan.isconnected():
        print('WiFi connected')
        print(wlan.ifconfig())
        return True

    print('Connecting WiFi {}'.format(WIFI_SSID))

    try:
        wlan.connect(WIFI_SSID, WIFI_PASSWORD)
    except OSError as e:
        print("WiFi HW Error: {}. Retrying...".format(e))
        machine.reset()

    for _ in range(30):
        if wlan.isconnected():
            print('WiFi connected')
            print(wlan.ifconfig())
            return True
        sleep(1)

    print('WiFi connection failed')
    return False


def connect_mqtt():
    global mqtt_client

    client_id = ubinascii.hexlify(machine.unique_id())
    mqtt_client = MQTTClient(
        client_id=client_id,
        server=MQTT_SERVER,
        port=MQTT_PORT,
        user=MQTT_USER,
        password=MQTT_PASSWORD
    )
    mqtt_client.set_callback(on_mqtt_message)
    mqtt_client.connect()
    mqtt_client.subscribe(MQTT_COMMAND_TOPIC)
    print('Connected to MQTT {}'.format(MQTT_SERVER))
    return mqtt_client


def publish_status(topic, event, ok, detail=''):
    payload = ujson.dumps({
        'node': NODE_ID,
        'event': event,
        'ok': ok,
        'detail': detail
    })
    mqtt_client.publish(topic, payload)


def set_screen(lines):
    if 'screen' not in globals() or 'oled' not in globals():
        return
    for index, line in enumerate(lines):
        screen[index] = line
    write_screen(oled, screen)


def send_key_update(command_packet):
    lora.println(command_packet)
    lora.receive()


def on_mqtt_message(topic, msg):
    global AES_KEY
    global HMAC_KEY
    global no_check_mode

    actual_topic = topic.decode() if type(topic) == bytes else topic
    config_topic = MQTT_COMMAND_TOPIC.decode() if type(MQTT_COMMAND_TOPIC) == bytes else MQTT_COMMAND_TOPIC

    if actual_topic != config_topic:
        return

    # Commande NOCHECK (activer/desactiver la verification HMAC)
    if msg.startswith(NOCHECK_PREFIX):
        ok, state = verify_nocheck_command(msg, HMAC_KEY)
        if ok:
            no_check_mode = state
            set_screen(['NO CHECK', 'ON' if state else 'OFF'])
            mqtt_client.publish(MQTT_DATA_TOPIC, b"nocheck_updated")
        return

    print("[MQTT] Ordre de Key Update recu.")
    ok, new_aes_key, new_hmac_key = verify_key_update(msg, HMAC_KEY)

    if ok:
        AES_KEY = new_aes_key
        HMAC_KEY = new_hmac_key
        mqtt_client.publish(MQTT_DATA_TOPIC, b"keys_updated")
        if 'lora' in globals():
            send_key_update(msg)
            lora.receive()
        set_screen(['KEYS OK', NODE_ID, 'updated'])
        print("[REKEY] Cles mises a jour avec succes et relayees.")
    else:
        mqtt_client.publish(MQTT_DATA_TOPIC, b"keys_rejected")
        set_screen(['KEYS BAD', NODE_ID])
        print("[REKEY] Echec : HMAC invalide.")
def publish_lora_message(node_id, plaintext, rssi, snr, secure):
    try:
        data_str = plaintext.decode() if type(plaintext) == bytes else str(plaintext)
    except Exception:
        data_str = str(plaintext)

    payload = ujson.dumps({
        'node': node_id,
        'data': data_str,
        'rssi': rssi,
        'snr': snr,
        'secure': secure
    })
    mqtt_client.publish(MQTT_DATA_TOPIC, payload)


def run():
    global lora
    global oled
    global screen

    if not connect_wifi():
        return

    connect_mqtt()
    lora = init_lora()
    oled, screen = init_oled()

    screen[0] = 'Passerelle'
    screen[1] = NODE_ID
    screen[2] = 'MQTT + LoRa'
    write_screen(oled, screen)

    while True:
        try:
            mqtt_client.check_msg()
        except OSError:
            connect_mqtt()

        if lora.received_packet():
            packet = lora.read_payload()
            rssi = lora.packet_rssi()
            snr = lora.packet_snr()
            mqtt_client.publish(MQTT_DATA_TOPIC, packet)

            if no_check_mode:
                screen[0] = 'RX RAW'
                screen[1] = packet[:16].decode() if type(packet) == bytes else packet[:16]
                screen[2] = 'RSSI {}'.format(rssi)
                screen[3] = 'SNR {}'.format(snr)
                screen[4] = 'NO-CHK'
            else:
                ok, node_id, plaintext = unpack_message(packet, AES_KEY, HMAC_KEY)
                if ok:
                    screen[0] = 'RX {}'.format(node_id)
                    screen[1] = plaintext.decode()[:16]
                    screen[2] = 'RSSI {}'.format(rssi)
                    screen[3] = 'SNR {}'.format(snr)
                    screen[4] = 'OK'
                else:
                    screen[0] = 'RX BAD'
                    screen[4] = 'HMAC ERR'
            write_screen(oled, screen)

            lora.receive()

        sleep(0.1)


if __name__ == '__main__':
    run()
\end{pycode}

\subsection{\texttt{init\_lora()} - une différence clé avec le capteur}
La configuration matérielle est identique à celle du capteur, mais \texttt{init\_lora()} appelle ici
\texttt{modem.receive()} immédiatement après la construction du pilote. La passerelle n'émet jamais de
mesure : son rôle est d'écouter le canal LoRa en continu, elle doit donc rester armée en réception dès
le démarrage.

\subsection{\texttt{connect\_wifi()} - connexion robuste}
Cette fonction va au-delà d'un simple \texttt{wlan.connect()} :
\begin{itemize}
    \item elle désactive puis réactive l'interface Wi-Fi (\texttt{wlan.active(False)} puis
    \texttt{True}) afin de repartir d'un état propre, utile après un \emph{soft-reset} de l'ESP32 ;
    \item elle court-circuite la tentative de connexion si le module est déjà connecté ;
    \item l'appel à \texttt{wlan.connect()} est protégé par un \texttt{try/except OSError} : certaines
    puces Wi-Fi ESP32 lèvent une erreur matérielle transitoire, à laquelle on répond par un
    \texttt{machine.reset()} complet plutôt que de rester bloqué ;
    \item une boucle de 30 tentatives, espacées d'une seconde, attend la connexion avant d'abandonner et
    de renvoyer \texttt{False}.
\end{itemize}

\subsection{\texttt{connect\_mqtt()} - identité et abonnement}
L'identifiant de connexion MQTT (\texttt{client\_id}) est dérivé de \texttt{machine.unique\_id()}, un
identifiant matériel propre à chaque puce ESP32, converti en chaîne hexadécimale. Cela garantit un
\texttt{client\_id} unique même si plusieurs passerelles partagent le même \texttt{config.py} : un broker
MQTT déconnecte en effet le premier client dès qu'un second se reconnecte avec le même identifiant.
La bibliothèque \texttt{umqttsimple} (un client MQTT minimaliste pour MicroPython) gère ensuite la
connexion TCP, l'authentification par utilisateur/mot de passe, puis l'abonnement immédiat au topic de
commande \texttt{MQTT\_COMMAND\_TOPIC}.

\subsection{\texttt{publish\_status()} et \texttt{set\_screen()} - fonctions utilitaires}
\texttt{publish\_status()} encapsule la publication d'un évènement de diagnostic au format JSON.
\texttt{set\_screen()} met à jour l'écran OLED de façon défensive : elle vérifie d'abord que les variables
globales \texttt{oled} et \texttt{screen} existent avant d'écrire, ce qui évite un plantage si la fonction
est appelée (par exemple depuis le callback MQTT) avant la fin de l'initialisation dans \texttt{run()}.

\subsection{\texttt{on\_mqtt\_message()} - traitement des commandes}
Cette fonction est enregistrée comme \emph{callback} et déclenchée par \texttt{umqttsimple} chaque fois
qu'un message arrive sur un topic souscrit. Elle gère deux types de commandes :

\paragraph{Commandes de mise à jour des clés (\texttt{KEYUPDATE})}
\begin{enumerate}
    \item elle compare le topic reçu au topic de commande attendu, en normalisant les deux côtés en
    chaînes de caractères (\texttt{topic} est parfois transmis en \texttt{bytes} par la bibliothèque MQTT) ;
    elle ignore silencieusement tout message qui ne concerne pas ce topic ;
    \item elle délègue la vérification cryptographique à \texttt{verify\_key\_update()} définie dans
    \texttt{security.py}, qui recalcule le HMAC de la commande avec la clé HMAC \textbf{actuelle} et la
    compare en temps constant ;
    \item en cas de succès : les nouvelles clés sont appliquées immédiatement (\texttt{global AES\_KEY},
    \texttt{global HMAC\_KEY}), un accusé \texttt{keys\_updated} est publié sur le topic de données, puis
    la commande est \textbf{relayée telle quelle vers le capteur par LoRa} via \texttt{send\_key\_update()} ;
    \item en cas d'échec, un accusé \texttt{keys\_rejected} est publié et les clés en cours
    \textbf{ne sont pas modifiées} : une commande invalide ne peut pas interrompre le service.
\end{enumerate}

\paragraph{Commandes de mode inspection (\texttt{NOCHECK})}
La passerelle reconnaît également les commandes préfixées par \texttt{NOCHECK\_PREFIX}. Leur HMAC est
d'abord vérifié par \texttt{verify\_nocheck\_command()} : si la signature est valide, la variable globale
\texttt{no\_check\_mode} est basculée à \texttt{True} ou \texttt{False} selon l'état demandé, et un
\texttt{nocheck\_updated} est publié sur le topic de données. En mode inspection activé, la boucle
\texttt{run()} affiche les trames reçues sans les vérifier (étiquette \texttt{NO-CHK} sur l'OLED), ce
qui permet d'observer le trafic brut pendant les tests sans modifier les clés.

\begin{explication}[Chaîne de confiance bout-en-bout]
Le client Linux signe la commande avec la clé HMAC qu'il croit être la clé courante. La passerelle la
vérifie avec \emph{sa propre} clé HMAC courante avant de l'appliquer, puis la relaie sans la modifier.
Le capteur effectue à son tour la même vérification avec sa propre clé courante. Une commande n'est donc
acceptée par un nœud que s'il partageait déjà la bonne clé HMAC au moment de la recevoir - ce qui
empêche un tiers ne connaissant pas la clé en cours de forcer une rotation.
\end{explication}

\subsection{\texttt{publish\_lora\_message()} - publication MQTT enrichie}
La fonction \texttt{publish\_lora\_message()} construit un message JSON structuré avant publication sur
le broker. Ce JSON regroupe l'identifiant du nœud émetteur, les données utiles déchiffrées, le niveau de
signal RSSI, le rapport signal/bruit SNR et un indicateur booléen \texttt{secure} précisant si le paquet
a bien passé la vérification HMAC. Cette structure facilite le traitement côté abonné MQTT : un client
peut filtrer immédiatement les messages sur le champ \texttt{secure}, et exploiter RSSI/SNR pour surveiller
la qualité du lien radio sans outillage supplémentaire.

\subsection{\texttt{run()} - boucle principale}
Après connexion Wi-Fi (avec retour anticipé si elle échoue), connexion MQTT puis initialisation LoRa et
OLED, la passerelle entre dans une boucle infinie qui, à chaque itération :
\begin{enumerate}
    \item appelle \texttt{mqtt\_client.check\_msg()}, qui interroge le socket MQTT de façon non
    bloquante et déclenche \texttt{on\_mqtt\_message()} si un message est en attente ; une
    \texttt{OSError} (connexion TCP perdue) entraîne une reconnexion immédiate via \texttt{connect\_mqtt()} ;
    \item vérifie si une trame LoRa est disponible (\texttt{lora.received\_packet()}) ; si oui, elle lit
    le paquet, le RSSI et le SNR, publie le paquet brut sur MQTT, puis tente de le déchiffrer pour
    actualiser l'écran OLED ;
    \item réarme systématiquement le module en réception avec \texttt{lora.receive()} - une étape
    \textbf{critique} : sans cet appel, le SX127x resterait dans l'état « réception terminée » et
    ignorerait toute trame suivante ;
    \item attend \SI{100}{\milli\second} (\texttt{sleep(0.1)}) avant de recommencer, un compromis entre
    réactivité et charge CPU.
\end{enumerate}

% =====================================================
\chapter{Client Python - supervision et sécurité}
\label{chap:client}

Le client Linux est structuré en \textbf{package Python} (\texttt{Projet/mqtt\_client/}) composé de
quatre fichiers aux responsabilités distinctes :
\begin{itemize}
    \item \texttt{config.py} : toutes les constantes (broker, topics, clés par défaut, préfixes de protocole) ;
    \item \texttt{security.py} : les primitives cryptographiques (AES-CBC, HMAC-SHA1, pack/unpack, rotation des clés, commande nocheck) ;
    \item \texttt{shell.py} : la classe \texttt{MqttShell} (interface interactive \texttt{cmd.Cmd}) et la fonction \texttt{print\_async()} ;
    \item \texttt{main.py} : le point d'entrée, la connexion MQTT, le traitement des messages et la boucle principale.
\end{itemize}

Cette séparation permet de réutiliser chaque module indépendamment et d'étendre le client sans toucher au
code cryptographique. Le client réimplémente côté CPython les mêmes primitives que \texttt{security.py}
côté MicroPython, afin que les messages chiffrés sur l'ESP32 puissent être vérifiés sur Linux et inversement.

\section{Constantes : \texttt{config.py}}

\begin{pycode}{Projet/mqtt\_client/config.py}
BLOCK_SIZE = 16
MESSAGE_PREFIX = b'INSARISF1'
KEYUPDATE_PREFIX = b'KEYUPDATE'
NOCHECK_PREFIX = b'NOCHECK'
DEFAULT_BROKER = 'serveur.iot.com'
DEFAULT_PORT = 1883
DEFAULT_USER = 'esp32'
DEFAULT_PASSWORD = 'esp32'
DEFAULT_DATA_TOPIC = 'LamBouba'
DEFAULT_COMMAND_TOPIC = 'echange/command/passerelle'
DEFAULT_AES_KEY = b'0123456789abcdef'
DEFAULT_HMAC_KEY = b'abcd'
\end{pycode}

Ce fichier centralise tous les paramètres modifiables en un seul endroit. Le préfixe \texttt{NOCHECK\_PREFIX}
identifie les commandes de mode inspection envoyées à la passerelle. Les constantes \texttt{MESSAGE\_PREFIX}
et \texttt{KEYUPDATE\_PREFIX} doivent être strictement identiques à celles du firmware, car elles font
partie intégrante du calcul HMAC.

\section{Primitives cryptographiques : \texttt{security.py}}

\begin{pycode}{Projet/mqtt\_client/security.py - bourrage et HMAC}
def pad(data):
    padding = BLOCK_SIZE - (len(data) % BLOCK_SIZE)
    return data + bytes([padding]) * padding

def unpad(data):
    padding = data[-1]
    if padding < 1 or padding > BLOCK_SIZE:
        raise ValueError('bad padding')
    return data[:-padding]

def hmac_sha1(key, message):
    return hmac.new(key, message, hashlib.sha1).digest()

def equal_bytes(a, b):
    return hmac.compare_digest(a, b)
\end{pycode}

\begin{itemize}
    \item \texttt{pad()}/\texttt{unpad()} implémentent un bourrage de type PKCS\#7 : on ajoute toujours
    entre 1 et 16 octets, chacun valant le nombre d'octets ajoutés, ce qui permet de retirer le bourrage
    sans ambiguïté en lisant le dernier octet ;
    \item \texttt{hmac\_sha1()} encapsule \texttt{hmac.new(...)} de la bibliothèque standard (RFC 2104) ;
    \item \texttt{equal\_bytes()} s'appuie sur \texttt{hmac.compare\_digest()}, une comparaison en
    \textbf{temps constant} pour résister aux attaques temporelles. Cette précaution est reprise
    manuellement dans le firmware, MicroPython ne fournissant pas \texttt{hmac.compare\_digest}.
\end{itemize}

\begin{pycode}{Projet/mqtt\_client/security.py - chiffrement AES-CBC}
def encrypt_aes_cbc(key, plaintext):
    nonce = os.urandom(BLOCK_SIZE)
    cipher = Cipher(algorithms.AES(key), modes.CBC(nonce), backend=default_backend())
    encryptor = cipher.encryptor()
    ciphertext = encryptor.update(pad(plaintext)) + encryptor.finalize()
    return nonce + ciphertext

def decrypt_aes_cbc(key, data):
    nonce = data[:BLOCK_SIZE]
    ciphertext = data[BLOCK_SIZE:]
    cipher = Cipher(algorithms.AES(key), modes.CBC(nonce), backend=default_backend())
    decryptor = cipher.decryptor()
    return unpad(decryptor.update(ciphertext) + decryptor.finalize())
\end{pycode}

Un IV (\texttt{nonce}) de 16 octets aléatoire est généré par \texttt{os.urandom()} à chaque chiffrement
et préfixé au texte chiffré. Deux messages identiques successifs produisent donc des textes chiffrés
différents, empêchant un observateur passif de repérer les répétitions dans le trafic.

\section{Empaquetage, messages et rotation des clés}

\begin{pycode}{Projet/mqtt\_client/security.py - pack/unpack}
def pack_message(node_id, plaintext, aes_key, hmac_key):
    encrypted = encrypt_aes_cbc(aes_key, plaintext)
    mac_data = MESSAGE_PREFIX + b'|' + node_id.encode() + b'|' + encrypted
    mac = hmac_sha1(hmac_key, mac_data)
    packet = node_id.encode() + b'|' + encrypted + b'|' + mac
    return base64.b64encode(packet).decode()

def unpack_message(packet, aes_key, hmac_key):
    try:
        raw = base64.b64decode(packet)
        node_id_bytes, rest = raw.split(b'|', 1)
        node_id = node_id_bytes.decode()
        encrypted, mac = rest.rsplit(b'|', 1)
        mac_data = MESSAGE_PREFIX + b'|' + node_id_bytes + b'|' + encrypted
        expected = hmac_sha1(hmac_key, mac_data)
        if not equal_bytes(mac, expected):
            return False, node_id, b''
        plaintext = decrypt_aes_cbc(aes_key, encrypted)
        return True, node_id, plaintext
    except Exception:
        return False, '', b''
\end{pycode}

\texttt{unpack\_message()} vérifie le HMAC \textbf{avant} de déchiffrer (approche \emph{verify-then-decrypt}).
Le \texttt{rsplit(b'|', 1)} pour extraire le HMAC (dernier champ) et le \texttt{split(b'|', 1)} pour
extraire l'identifiant (premier champ) tolèrent un \texttt{node\_id} contenant accidentellement le séparateur.

\begin{validation}[Vérification HMAC avant déchiffrement]
Vérifier le HMAC avant de déchiffrer empêche un oracle de padding : si l'on déchiffrait d'abord, un
message d'erreur de bourrage pourrait renseigner un attaquant sur le contenu du texte chiffré. Seuls les
paquets signés avec la bonne clé HMAC sont déchiffrés, sur toute la chaîne.
\end{validation}

\begin{pycode}{Projet/mqtt\_client/security.py - rotation et nocheck}
def build_key_update(new_aes_key, new_hmac_key, current_hmac_key):
    nonce = os.urandom(8)
    command = KEYUPDATE_PREFIX + b'|' + nonce + b'|' + new_aes_key.hex().encode() + b'|' + new_hmac_key.hex().encode()
    mac = hmac_sha1(current_hmac_key, command)
    return command + b'|' + mac.hex().encode()

def build_nocheck_command(state_bool, current_hmac_key):
    nonce = os.urandom(8)
    state = b'1' if state_bool else b'0'
    command = NOCHECK_PREFIX + b'|' + nonce + b'|' + state
    mac = hmac_sha1(current_hmac_key, command)
    return command + b'|' + mac.hex().encode()
\end{pycode}

Le format de la commande de rotation est :
\texttt{KEYUPDATE | nonce(8 octets) | nouvelle\_clé\_AES(hex) | nouvelle\_clé\_HMAC(hex) | HMAC(commande)}.
Le HMAC est calculé avec la clé HMAC \textbf{courante}, ce qui garantit que seul le détenteur des secrets
actuels peut autoriser une rotation.

\texttt{build\_nocheck\_command()} construit la commande de bascule du mode inspection, également signée
avec la clé HMAC courante. Un nonce de 8 octets garantit que deux commandes successives vers le même état
produisent des paquets différents.

\begin{explication}[Rôle du nonce dans les commandes]
Le nonce aléatoire de 8 octets empêche un observateur de reconnaître une commande à sa simple structure
binaire, et rend impossible la réutilisation d'une commande interceptée.
\end{explication}

\section{Interface interactive : \texttt{shell.py}}

\begin{pycode}{Projet/mqtt\_client/shell.py - MqttShell (extrait)}
class MqttShell(cmd.Cmd):
    intro = "Client MQTT/LoRa securise interactif."
    prompt = "(mqtt) "

    def do_listen(self, arg):
        self.client.subscribe(self.args.data_topic)
        self.listening = True

    def do_update_keys(self, arg):
        # Publie la commande de rotation et active le mode double-decodage
        new_aes_key, new_hmac_key = [a.encode() for a in arg.split()]
        command = build_key_update(new_aes_key, new_hmac_key, self.args.hmac_key)
        self.client.publish(self.args.command_topic, command)
        self.args.pending_update = True
        self.args.fallback_keys = (self.args.aes_key, self.args.hmac_key)
        self.args.aes_key = new_aes_key
        self.args.hmac_key = new_hmac_key

    def do_nocheck(self, arg):
        state = (arg.strip().lower() == 'on')
        command = build_nocheck_command(state, self.args.hmac_key)
        self.client.publish(self.args.command_topic, command)

    def do_status(self, arg):
        print(f"Broker: {self.args.broker}:{self.args.port}")
        print(f"Cle AES actuelle: {self.args.aes_key.decode()}")
        print(f"Cle HMAC actuelle: {self.args.hmac_key.decode()}")
\end{pycode}

La classe \texttt{MqttShell} hérite de \texttt{cmd.Cmd}, le module de la bibliothèque standard Python
pour les interfaces en ligne de commande interactives. Chaque méthode préfixée \texttt{do\_} définit
automatiquement une commande. La boucle \texttt{cmdloop()} gère la lecture, la complétion par tabulation
et l'aide automatique depuis les docstrings. Les commandes disponibles sont :

\begin{itemize}
    \item \texttt{listen} / \texttt{stop\_listen} : s'abonner ou se désabonner du topic de données ;
    \item \texttt{update\_keys <aes> <hmac>} : lancer une rotation de clés et activer le mode double-décodage ;
    \item \texttt{nocheck on|off} : activer ou désactiver la vérification HMAC sur la passerelle ;
    \item \texttt{publish <node\_id> <message>} : publier un message de test sécurisé ;
    \item \texttt{status} : afficher broker, topic et clés actives ;
    \item \texttt{quit} / \texttt{exit} : fermer le client.
\end{itemize}

La fonction \texttt{print\_async()} assure un affichage non-intrusif des messages MQTT reçus en arrière-plan :
elle efface la ligne courante, affiche le message, puis réaffiche le prompt et le contenu en cours de saisie
via \texttt{readline.get\_line\_buffer()}, de sorte que les notifications entrantes n'interrompent pas la
frappe d'une commande.

\section{Point d'entrée : \texttt{main.py}}

\begin{pycode}{Projet/mqtt\_client/main.py - on\_message avec double-decodage}
def on_message(client, userdata, message, args, shell=None):
    raw_payload = message.payload.decode()
    SYSTEM_MESSAGES = ("keys_updated", "keys_rejected", "nocheck_updated")
    if raw_payload in SYSTEM_MESSAGES:
        print_async(f'\n[SYSTEM] Passerelle: {raw_payload}', shell)
        return

    pending = getattr(args, 'pending_update', False)
    valid, node_id, plaintext = unpack_message(message.payload, args.aes_key, args.hmac_key)

    if pending and not valid:
        # Essai avec les anciennes cles (Trial Decryption)
        fb_aes, fb_hmac = args.fallback_keys
        fb_valid, fb_node_id, fb_plaintext = unpack_message(message.payload, fb_aes, fb_hmac)
        if fb_valid:
            args.message_count += 1
            data = fb_plaintext.decode()
            if data == "KEYS_FAILED" or args.message_count >= 2:
                # Rollback : restaurer les anciennes cles localement et republier
                args.aes_key, args.hmac_key = fb_aes, fb_hmac
                args.pending_update = False
                command = build_key_update(fb_aes, fb_hmac, args.active_keys[1])
                client.publish(args.command_topic, command)
    elif pending and valid:
        if plaintext.decode() == "KEYS_OK":
            args.pending_update = False  # Rotation confirmee
    elif valid:
        print_async(f'\n[OK] {node_id} {plaintext.decode()}', shell)
    else:
        print_async(f'\n[BAD] message rejected\nraw: {raw_payload}', shell)
\end{pycode}

\subsection{\texttt{on\_message()} - traitement et rollback automatique}
Cette fonction est le cœur du traitement des messages MQTT côté client :
\begin{enumerate}
    \item les messages système de la passerelle (\texttt{keys\_updated}, \texttt{keys\_rejected},
    \texttt{nocheck\_updated}) sont affichés directement sans tentative de déchiffrement ;
    \item hors rotation en attente, \texttt{unpack\_message()} est appelé avec les clés courantes ;
    \item en mode \texttt{pending\_update} (rotation demandée mais pas encore confirmée), un
    \textbf{double-décodage} est appliqué : on essaie d'abord les nouvelles clés, puis les anciennes
    (\texttt{fallback\_keys}) si les nouvelles échouent ;
    \item si deux messages consécutifs sont déchiffrables uniquement avec les anciennes clés, ou si le
    capteur répond \texttt{KEYS\_FAILED}, le client déclenche un \textbf{rollback automatique} : il
    restaure ses clés localement et publie une nouvelle commande de rotation vers les anciennes valeurs,
    à destination de la passerelle ;
    \item si le capteur répond \texttt{KEYS\_OK} déchiffrable avec les nouvelles clés, la rotation est
    confirmée et \texttt{pending\_update} est désactivé.
\end{enumerate}

\begin{pycode}{Projet/mqtt\_client/main.py - boucle principale}
def main():
    args = parser.parse_args()
    client = make_mqtt_client(args)
    shell = MqttShell(args, client)
    client.on_message = lambda c, u, m: on_message(c, u, m, args, shell)
    client.loop_start()   # Thread MQTT en arriere-plan
    shell.cmdloop()       # Boucle interactive au premier plan
    client.loop_stop()
    client.disconnect()
\end{pycode}

\subsection{Architecture asynchrone}
La boucle MQTT est lancée en \textbf{arrière-plan} via \texttt{client.loop\_start()}, qui crée un thread
dédié. Le fil principal est ainsi libre pour la boucle interactive \texttt{shell.cmdloop()}. Les deux
boucles s'exécutent en parallèle : le shell attend la saisie pendant que le thread MQTT reçoit et traite
les messages, en appelant \texttt{print\_async()} pour afficher les notifications sans perturber le prompt.

\begin{explication}[Format des clés sur la ligne de commande]
\texttt{parse\_bytes()} convertit un argument de ligne de commande directement en octets ASCII
(\texttt{value.encode()}) : taper \texttt{--aes-key 0123456789abcdef} produit les 16 octets de la clé
attendue, cohérents avec \texttt{DEFAULT\_AES\_KEY = b'0123456789abcdef'} dans \texttt{config.py}.
Pour les nouvelles clés transmises par \texttt{update\_keys}, \texttt{build\_key\_update()} les encode en
hexadécimal (\texttt{new\_aes\_key.hex()}) pour les transporter sans ambiguïté de séparateur.
\end{explication}


% =====================================================
\chapter{Sécurisation des échanges}
\label{chap:securite}

\section{Chiffrement AES des données}
Le chiffrement est réalisé en AES-CBC avec une clé de 16 octets. Chaque paquet contient un nonce de
16 octets utilisé comme IV AES. Ce nonce est aléatoire, ce qui évite que deux messages identiques
produisent le même texte chiffré.

\section{Authentification et intégrité des messages}
L'intégrité et l'origine du message sont garanties par un HMAC-SHA1. Le HMAC est calculé sur le préfixe
\texttt{INSARISF1}, l'identifiant du nœud et la partie chiffrée. Cette construction empêche qu'un
attaquant modifie le message sans connaître la clé HMAC.

\section{Protocole de mise à jour sécurisée des clés}
Le client Linux peut publier une commande de mise à jour des clés sur le topic de commande. Cette
commande contient une nouvelle clé AES, une nouvelle clé HMAC et un HMAC calculé avec la clé HMAC
actuelle. La passerelle vérifie la commande, met à jour ses clés puis la transmet au capteur par LoRa.
Le capteur vérifie à son tour la commande avant d'appliquer les nouvelles clés.

\section{Analyse de sécurité}
\begin{validation}[Ce que le protocole protège]
Le chiffrement AES-CBC garantit la confidentialité face à une écoute passive du canal radio ou MQTT, et
le HMAC-SHA1 garantit qu'un message modifié ou forgé sans connaître la clé sera détecté et rejeté.
\end{validation}

\begin{validation}[Vérification HMAC avant déchiffrement]
Le protocole vérifie systématiquement le HMAC \textbf{avant} de déchiffrer : cette approche
\emph{verify-then-decrypt} évite qu'un attaquant exploite les erreurs de déchiffrement pour en déduire
des informations sur le contenu (oracle de padding). Seuls les paquets signés avec la bonne clé
HMAC sont déchiffrés, sur l'ensemble de la chaîne : radio LoRa, broker MQTT, et client Linux.
\end{validation}

% =====================================================
\chapter{Résultats expérimentaux}
\label{chap:resultats}

\section{Mise en service de la chaîne complète}
Le client Linux interactif (\texttt{mqtt\_client/main.py}) a été lancé en parallèle des deux firmwares
ESP32, chacun affichant son propre journal sur sa liaison série dans un panneau de terminal dédié. La
figure~\ref{fig:capture-demarrage} montre cet environnement de test : à gauche le journal du capteur
(émission \texttt{[TX]} d'un paquet puis ouverture de la fenêtre d'écoute des clés), à droite celui de la
passerelle (réception \texttt{[RX]}, niveau RSSI et SNR), et en haut le client Python, qui liste au
démarrage l'ensemble des commandes disponibles : \texttt{listen}, \texttt{update\_keys}, \texttt{nocheck},
\texttt{publish}, \texttt{status}, \texttt{stop\_listen}.

\section{Réception et déchiffrement des messages sécurisés}
La commande \texttt{listen} place le client en écoute sur le topic \texttt{LamBouba}. Chaque trame brute
reçue (\texttt{MQTT raw}) est vérifiée puis déchiffrée : en cas de succès, le client affiche
\texttt{[OK] capteur1} suivi du contenu JSON en clair (température, batterie, compteur). La
figure~\ref{fig:capture-listen} illustre cette réception pour deux messages consécutifs, avec une
température croissante d'une mesure à l'autre, cohérente avec la rampe simulée par \texttt{build\_payload()}.

\section{Mise à jour des clés}
La commande \texttt{update\_keys} publie une commande de rotation sur le topic
\texttt{echange/command/passerelle}. La passerelle applique les nouvelles clés puis relaie la commande au
capteur par LoRa, qui confirme à son tour la mise à jour. Le client affiche alors
\texttt{[SYSTEM] Passerelle: keys\_updated} puis \texttt{[CONFIRMED] Mise à jour des clés validée par le
capteur}, comme le montre la figure~\ref{fig:capture-updatekeys}. La commande \texttt{status} permet
ensuite de relire les clés AES et HMAC actuellement actives sur la passerelle.

\section{Rejet d'un message altéré}
Lorsqu'un message ne passe pas la vérification HMAC, le client l'identifie avec
\texttt{[BAD] message rejected} et affiche la trame brute reçue, sans que cela n'interrompe la réception
des messages valides suivants : la figure~\ref{fig:capture-rejet} montre un message rejeté entouré de deux
messages acceptés normalement, illustrant que la vérification s'effectue indépendamment pour chaque trame.

\section{Désactivation de la vérification}
La commande \texttt{nocheck on} permet de désactiver volontairement la vérification HMAC côté passerelle,
afin d'observer la réception des trames en clair (étiquetées \texttt{NO-CHK} sur l'écran de la passerelle
et publiées en MQTT sous forme \texttt{RX RAW}). La figure~\ref{fig:capture-nocheck} illustre ce mode de
diagnostic, utile pour inspecter le contenu brut d'un paquet pendant les tests.

\section{Rotation avec mauvaises clés et rollback automatique}
La figure~\ref{fig:capture-rollback} illustre le comportement du système lorsqu'une rotation échoue côté
capteur : la passerelle relaie la nouvelle clé proposée, mais le capteur la rejette
(\texttt{[BAD\_KEYS] (Old Keys Valid)}) et continue de signer ses messages avec son ancienne clé HMAC. Le
client journalise un \texttt{[ROLLBACK]} et la passerelle revient à son tour aux clés précédentes, ce qui
rétablit immédiatement la cohérence entre les trois composants sans interrompre le flux de mesures.

\section{Autres outils de test disponibles}
En complément du mode interactif, le projet fournit des outils en ligne de commande pour des tests rapides
ou scriptés. Un test direct du broker, indépendant du chiffrement, peut être réalisé avec
\texttt{mosquitto\_pub} :
\begin{shellcode}{test de publication MQTT}
mosquitto_pub -h serveur.iot.com -u esp32 -P esp32 -t LamBouba -m coucou
\end{shellcode}
La sous-commande \texttt{publish-test} de \texttt{mqtt\_client/main.py} permet de publier un unique paquet
sécurisé sans passer par le capteur LoRa :
\begin{shellcode}{publish-test}
python3 -m mqtt_client --broker serveur.iot.com \
  --mqtt-user esp32 --mqtt-password esp32 \
  --aes-key 0123456789abcdef --hmac-key abcd \
  publish capteur1 coucou
\end{shellcode}

\section{Synthèse}
L'ensemble de ces tests couvre le cycle de vie complet du protocole : émission et réception normales,
rotation de clés réussie, rejet ciblé d'un message altéré au sein d'un flux par ailleurs valide, mode de
diagnostic sans vérification, et reprise automatique après une rotation de clés défaillante. Le système se
comporte de façon cohérente sur l'ensemble de ces scénarios.

% =====================================================
\chapter{Conclusion}
\label{chap:conclusion}

Le projet permet de mettre en œuvre une chaîne complète IoT LoRa/MQTT avec vérification cryptographique
des messages. Le capteur produit des mesures, la passerelle les transporte vers MQTT et le client Linux
vérifie les messages reçus. Les mécanismes AES-CBC et HMAC-SHA1 assurent respectivement la confidentialité
et l'intégrité des échanges.

Les améliorations principales seraient l'ajout d'un compteur anti-rejeu, le remplacement d'AES-CBC par
AES-GCM et l'ajout d'une journalisation plus détaillée côté passerelle. Le script Bash de publication
permet désormais de générer et d'envoyer automatiquement un flot de paquets sécurisés avec des messages
différents, ce qui facilite les tests de réception et de robustesse.

% =====================================================
\appendix
\chapter{Annexes}

\section{Code source complet}
Les fichiers principaux du projet sont :
\begin{itemize}
    \item \texttt{Projet/firmware/capteur/main.py}, \texttt{config.py}, \texttt{security.py} ;
    \item \texttt{Projet/firmware/passerelle/main.py}, \texttt{config.py}, \texttt{security.py} ;
    \item \texttt{Projet/mqtt\_client/main.py}, \texttt{config.py}, \texttt{security.py}, \texttt{shell.py}.
\end{itemize}

\section{Schémas de câblage}
Les ESP32 sont câblés aux modules LoRa SX127x en SPI. Les broches utilisées sont définies dans les
fichiers \texttt{config.py} du capteur et de la passerelle. Les deux modules doivent partager la masse,
l'alimentation, SCK, MOSI, MISO, NSS, RST et DIO0.

\section{Captures d'écran}

\begin{figure}[H]
    \centering
    \includegraphics[width=\textwidth]{captures/demo1_demarrage_tx_rx.png}
    \caption{Démarrage du client interactif et premier échange capteur~$\rightarrow$~passerelle (terminaux MicroPython à gauche et à droite, client Python en haut)}
    \label{fig:capture-demarrage}
\end{figure}

\begin{figure}[H]
    \centering
    \includegraphics[width=\textwidth]{captures/demo2_listen_securise.png}
    \caption{Mode \texttt{listen} - réception, vérification HMAC et déchiffrement de messages sécurisés (\texttt{[OK] capteur1 \{...\}})}
    \label{fig:capture-listen}
\end{figure}

\begin{figure}[H]
    \centering
    \includegraphics[width=\textwidth]{captures/demo3_update_keys_succes.png}
    \caption{Commande \texttt{update\_keys} - rotation des clés relayée à la passerelle puis confirmée par le capteur (\texttt{[CONFIRMED]})}
    \label{fig:capture-updatekeys}
\end{figure}

\begin{figure}[H]
    \centering
    \includegraphics[width=\textwidth]{captures/demo4_message_rejete.png}
    \caption{Rejet d'un message altéré au sein du flux (\texttt{[BAD] message rejected}), les messages suivants correctement signés continuant d'être acceptés}
    \label{fig:capture-rejet}
\end{figure}

\begin{figure}[H]
    \centering
    \includegraphics[width=\textwidth]{captures/demo5_nocheck_on.png}
    \caption{Commande \texttt{nocheck on} - désactivation volontaire de la vérification HMAC côté passerelle pour observer la réception brute (\texttt{NO-CHK})}
    \label{fig:capture-nocheck}
\end{figure}

\begin{figure}[H]
    \centering
    \includegraphics[width=\textwidth]{captures/demo6_rollback_status.png}
    \caption{Rotation de clés avec un mauvais HMAC (\texttt{[BAD\_KEYS]}), \texttt{[ROLLBACK]} automatique vers les clés précédentes, puis commande \texttt{status} récapitulant l'état courant de la passerelle}
    \label{fig:capture-rollback}
\end{figure}

\end{document}
